\section{The call-by-name $\lambda$-calculus}
\label{s:cbnNEW}


The relationship between OGS and $\pi$-calculus representations of the 
$\lambda$-terms examined in Sections~\ref{s:encoGames} and \ref{s:cogs} is not specific to
call-by-value.  A similar relationship can be established for call-by-name. 
We summarise here the main aspects and the main results. 
Further details may be found in Appendix~\ref{a:cbn}.


We maintain terminologies and notations in the call-by-value sections. 
The call-by-name language is given in Figure~\ref{f:cbnG}.
Its reduction $\redn$ is defined by the following two rules:
 
 \begin{mathpar}
   \inferrule*{ }{(\lambda x.M) N \redn M\subst{x}{N}} \qquad
   \inferrule*{M \redn N}{E[M] \redn E[N]}
 \end{mathpar}
 
%\subsection{CBN Operational Game Semantics}

The OGS for call-by-name uses
three kinds of names: 
\emph{continuation names} (ranged over by $p,q$), 
\emph{variable names} (ranged over by $x,y$),
and 
\emph{value names} (ranged over by $v,w$). 
Continuation names are used as in call-by-value; 
variable names act as pointers to 
$\lambda$-terms and correspond to the variables of the $\lambda$-calculus;
value names are pointers to values.

The OGS uses the syntactic categories in Figure~\ref{f:cbnG}, 
where $L$ range over extended $\lambda$-terms, 
which are needed to accommodate value names. 
We write $\LaoP$ for the set of extended $\lambda$-terms.
\begin{figure*}
\[ 
\begin{array}{rclr}
E & := &  EM \midd \contexthole  &  \mbox{evaluation contexts}\\
V & := &  \lambda x. M \midd v &  \mbox{values} \\
\gamma & := & 
\pmap{v}{V}\cdot\gamma' \midd 
\pmap{q}{(E,p)}\cdot\gamma' \midd 
\pmap{x}{M}\cdot\gamma' \midd 
\emptymap  &  \mbox{environments} \\
F & := & \conf{\LL,p,\gamma ,\phi } \midd 
 \conf{\gamma,\phi} \midd \initconf{M}{\phi} & \mbox{configurations} \\
\LL & := &  M \midd E [v M] \midd v   & \mbox{extended $\lambda$-terms}\\
\end{array} 
\] 
\caption{Grammars for the call-by-name language and \AOGS configurations}
\label{f:cbnG}
\end{figure*} 

\dsOLD{revise Figure~\ref{f:cbnG}}

\noindent
In a call-by-name setting, Player Questions are of two possible shapes:
 \begin{itemize}
  \item \emph{Player Value-Question} (PVQ) $\questPV{v}{x}{p}$,
  receiving a variable $x$ and a continuation name $p$ through a value name $v$;
  \item \emph{Player Term-Question} (PTQ) $\questPT{x}{p}$,
  receiving a continuation name $p$ through a variable $x$.
 \end{itemize}
\noindent
 Opponent Questions are also of two different shapes too:
  \begin{itemize}
   \item \emph{Opponent Value-Question} (OVQ) $\questOV{v}{x}{p}$,
  sending a variable $x$ and a continuation name $p$ through a value name $v$;
  \item \emph{Opponent Term-Question} (OTQ) $\questOT{x}{p}$,
  sending a continuation name $p$ through a variable $x$.
%   \item \emph{Initial Opponent Question} (IOQ) $\questOinit{p}$ 
%   that sends %a sequence of variables $\tilde{x}$ and 
%   a continuation name $p$ 
%   through the initial continuation name $\initName$.  
 \end{itemize}
  

 
 
 \begin{figure*}
 \setlength{\arraycolsep}{3pt}%
 \[
  \begin{array}{l|lll@{}l}
   (\tau) & \conf{M,p,\gamma,\phi} & \arrA{\ \tau \ } & 
     \conf{N,p,\gamma,\phi} & \text{ when } M \redn N\\
%
   (PA) & \conf{V,p,\gamma,\phi} & \arrA{\ansP{p}{v}} & 
     \conf{\gamma \cdot \pmap{v}{V},\phi \uplus \{v\}} \\
%
   (PVQ) & \conf{E[vM],p,\gamma,\phi} & \arrA{\questPV{v}{y}{q}} & 
     \conf{\gamma \cdot \pmap{y}{M}\cdot \pmap{q}{(E,p)},\phi \uplus \{y,q\}} \\
%
   (PTQ) & \conf{E[x],p,\gamma,\phi} & \arrA{\questPT{x}{q}} & 
     \conf{\gamma \cdot \pmap{q}{(E,p)},\phi \uplus \{q\}} \\
%
   (OA) & \conf{\gamma\cdot \pmap{p}{(E,q)},\phi} & \arrA{\ansO{p}{v}} & 
     \conf{E[v],q,\gamma,\phi \uplus \{v\}} \\ %& \text{ when } \gamma(p) = (E,q)  \\
%
   (OVQ) & \conf{\gamma \cdot [v \mapsto V],\phi} & \arrA{\questOV{v}{y}{p}} & 
     \conf{V y,p,\gamma,\phi \uplus \{p,y\}} & \text{ when } \gamma(v) = V\\
%
   (OTQ) & \conf{\gamma,\phi} & \arrA{\questOT{x}{p}} & 
     \conf{M,p,\gamma,\phi \uplus \{p\}} & \text{ when } \gamma(x) = M\\
%
   (IOQ) & \initconf{\phi}{M} & \arrA{\questOinit{p}} &
     \conf{M,p,\emptymap,\phi \uplus \{p\}}
%    (IOQ) & \conf{M} & \ired{\questOinit{\tilde{x}}{p}} &
%      \conf{M,p,\emptymap,\{\tilde{x},p\}} & \text{ with } \FV(M) = \{\tilde{x}\}
%
  \end{array}
 \]
 \caption{CBN OGS LTS}
\label{fig:scbn-lts}
\end{figure*}


%  \begin{theorem}
%   The sequence of actions generated by an initial configuration $\initconf{\phi}{v}{M}$ is a justified trace.
%  \end{theorem}
% 
% 
% \finish{what is the theorem above for?}



\begin{figure}
\Mybar 
{\small
Encoding of $\Lao$ into  \piI processes: 
\[\begin{array}{rcl}
\encoN{ \lambda  x. M}   & \defi & 
\bind p
  {  \bout p v  .   \inp v {x,q}  .
\encoNa M q}
 \\[\mypt]
\encoN{  x}   & \defi & \bind p
 { \bout xr . \link r p   } \\[\mypt]
\encoN{M N } & \defi &  \bind p
 {\res q \Big(
\encoNa{M}{ q} |
 \inp q v . { \bout v { x, {p'}} . ( \link {p'} p |  ! \inp x r. {\encoNa{ N }{ r}})  }
\Big)
}
\end{array}
   \]
$ $ \dotfill $ $ 

Encoding of    $\LaoP$:
\[ 
\begin{array}{rcl}
\encoN {E [v M]} & \defi & \bind p \bout v {x,r}. ( \encoN { \pmap r{(E,p)}} | ! \inp x
                       {q}. \encoNa Mq   ) 
\\[\mypt]
\encoN {v } & \defi & \bind p \bout p w . \link wv
\end{array} 
\]
$ $ \dotfill $ $ 

Encoding of  OGS environments:
\[ 
\begin{array}{rcl}
\encoN { \pmap{v}{V} \cdot \gamma   } & \defi & 
\left\{ \begin{array}{ll} 
 \inp v {x,q}. \encoNa M q  | \encoN{ \gamma}
&      %  ~\quad 
 \mbox{if $V  = \lambda x. M$ } \\[\mypt]
\link v w  | \encoN{ \gamma}
   & %    ~\quad  
\mbox{if $V  =  w$ } 
\end{array}  \right. 
\\[\myptt]
\encoN { \pmap{x}{M}  \cdot \gamma } & \defi &  ! \inp x q . \encoNa M q   | \encoN{ \gamma} \\[\mypt]
\encoN {\pmap{q}{(E,p)}  \cdot \gamma  } & \defi & 
\left\{ \begin{array}{@{}l@{\;}l@{}} 
 q(v).\bout v {x,r}.(\encoN{\pmap{r}{E',p}} | !\inp x q.\encoNa Mq)
  | \encoN{\gamma}
& \mbox{if $E=E'[\contexthole\,M]$} \\[\mypt]
\link qp   | \encoN{ \gamma} & \mbox{if $E=\contexthole$} 
\end{array}  \right. 
\\[\myptt]
\ifcutCBNLMCS
\multicolumn{3}{l}{ \mbox{(An alternative, but less confortable for
    proofs)} }
\\[\myptt]\multicolumn{3}{l}{ \mbox{(above: M is the first argument, below it is the last)} }
\\[\myptt]
\encoN {\pmap{q}{(E,p)}  \cdot \gamma  } & \defi & 
\left\{ \begin{array}{ll} 
\res r ( \pmap{q}{E',r }
|
 r(v). \bout v {x,p'}. 

(\link {p'}p |
 ! \inp x {q'} . \encoNa M{q'} )
  | \encoN{ \gamma}
&      %  ~\quad 
 \mbox{if $E   = 
E' M
$ } \\[\mypt]
\link qp   | \encoN{ \gamma} &      %  ~\quad 
 \mbox{if $E   = \contexthole$ } 
\end{array}  \right. 
\\[\myptt]
\fi
\encoN { \emptyset }  & \defi & \nil
\end{array} 
\]
$ $ \dotfill $ $ 

Encoding of  OGS configurations: 
\[ 
\begin{array}{rcl}
\encoN {\conf {M,p,\gamma , \phi}} & \defi & \pairP { \encoNa M p | \encoN \gamma }
\\[\mypt]
\encoN {\conf {v,p,\gamma , \phi}} & \defi & \pairP { \encoNa v p | \encoN \gamma }
\\[\mypt]
\encoN {\conf {E[v M],p,\gamma , \phi}} & \defi & \pairP { \encoNa {E[vM]} p | \encoN \gamma
                                             }
\\[\mypt]
\encoN {\conf {\gamma , \phi}} & \defi & \pairP { \encoN \gamma }
\end{array} 
\]
}% end small

\caption{The encoding of
call-by-name $\lambda$-calculus and OGS  into \piI }
\label{f:namePI}
\Mybar 
\end{figure} 



\noindent
The encodings of call-by-name 
$\lambda$-calculus,  OGS environments and configurations, and associated  syntactic categories,
 are reported in Figure~\ref{f:namePI}.
The \piI representation uses the same 3 types of  names as the OGS representation. 

As for call-by-value, so in call-by-name the encoding makes use of 
 \emph{link} processes, which are linear structures for continuation
 names and replicated structures for the other names.  We recall their
 definition. 
% When using \piI, a common recursive structure are the \emph{link}
% processes.  In  a  link from $c$ to $b$, written $\link  cb$, 
% names $b$ and $c$ are ``connected'', in the
% sense that  a process performing an  output at $c$ and a process
% performing an input at $b$ can interact, asynchronously,  through the
% link.  In other words, 
% a link behaves a little like  a name buffer: It receives names
% at one end-point and transmit names  at the other end-point. However, the
% latter names are not the same as the former names --- as it  would be
% in a  real buffer --- but, instead,   are {\em linked\/} to them:
% This accounts 
% for the recursion in the definition of links.
For tuples of names 
 $\til u = u_1, \ldots , u_n$
 and
$\til v =  v_1, \ldots , v_n$   
we write $\link{\til u}{\til v}$
to 
  abbreviate 
   $\link{u_1}{v_1} | \ldots |  \link{u_n}{v_n}$:
% The
% definition of  a link $\link{}{}$  (for call-by-name) is: 
$$
\begin{array}{rcl}
\link{}{} &\defi &
\left\{ \begin{array}{l}
 \abs{x,y} ! \inp x {p}.\bout {y} {q}
           . {\fwd{ q}{ p}}\\[\myptSmall] 
 \multicolumn{1}{l}{  
        ~\quad  \mbox{if $x,y$ are variable names}
% \abs{v,w}  \inp v {x,p}.\bout {w} {y,q}
%           . ({\link{ y}{ x}}|  {\link{ q}{ p}} )\\
% \multicolumn{1}{r}{  
%        ~\quad  \mbox{if $v,w$ are value variable names}
                   }              \\[\mypt]
\abs{a,b}   \inp a { \til c }.\bout b {\til d}
          .{\link{ \til d}{ \til c}})  \\
\multicolumn{1}{r}{  
          ~\quad \mbox{otherwise}
}\end{array} \right. 
\end{array}
 $$


\noindent
We have:
\[ 
\begin{array}{rcl}
\encoN {E [v ]} & = & 
\left\{ \begin{array}{ll}
\encoN {E' [vM ]}  & 
\mbox{if $E =
E'[\:  \contexthole M   \: ] %E' M
$}
\\[\mypt]
\encoN {v}  & 
\mbox{if $E = \contexthole$}
\end{array} \right. 
\end{array} 
\]


\dsJ{I have rectified above: 
{if $E =
E'[\:  \contexthole M   \: ] %E' M
$} rather than 
{if $E =
E' M
$}
% The above seems wrong to me now. 
% It seems to me that, if $E = [.]M \til M$ then 
% we have  $E[v] =  E'[vM] $ for $E' = [.]\til M$.
% But maybe the thing above is not used
}


\begin{lemma}
\label{l:Ex}
$ 
\begin{array}{rcl}
\encoN {E [x]} & = & \bind p 
\bout x {q}.  \encoN { \pmap q{(E,p)}} 
\end{array} 
$
\end{lemma} 
\noindent
The results of operational correspondence between OGS and \piI under the \opLTS are the
same as for call-by-value. 
We only report the operational correspondence for strong transitions and the 
analogue of Corollary~\ref{c:fa_traces_bis}.  


\begin{theorem}%[from OGS to \piI, strong transitions]
\label{t:opcorrCON_cbnshort}
$ $ \begin{enumerate}

% \item   If $F \arr{\initName (p)} F'$ then 
% $\encoN{F } \arr  {\initName (p)}  
%  \encoN{F'}$ ; 

\item  If $F \longrightarrowA F'$ then 
$ \encoN F \longrightarrowPI \;   \contr \encoN {F'}$; 

\item  If $F \arrA{\ell} F'$  %where $\mu$ is an input or an output action
% not $\tau$ or of the form 
% $\initName (p)$ 
then 
$ \encoN F \arr{\mu}  \; \wbPI  \encoN  {F'}$.
\end{enumerate}
\end{theorem} 


\begin{theorem} %[from \piI to  OGS, sequential LTS]
\label{t:opcorrCON3cbn}
$ $ \begin{enumerate}
% \item  If
% $\encoConV{\FF } \arrN  {\initName (p)} \PP$, then there is $\FF'$ 
% such that
%  $\FF \arr{\initName (p)} \FF'$ 
% and $\PP = \encoConV{\FF' }$ ;

\item If 
$ \encoN \FF \arrN \tau \PP$ then there is $\FF'$ such that 
$\FF \longrightarrowA \FF'$  and  
$\PP   \contr \encoN {\FF'}$; 


\item If
$ \encoN \FF \arrN\ell \PP$
 then there is $\FF'$ such that
 $\FF \arrA\ell \FF'$
and $\PP \wbPI \encoN {\FF'}$.
\end{enumerate} 
\end{theorem} 


\begin{corollary}
\label{c:CBNtraces}
For any configuration $\FF$ and trace $s$,
 we have 
  $\FF \ArrA s$   iff 
 $ \encoN \FF  \ArrN s$.
\end{corollary} 

\subsection{Concurrent Operational Game Semantics}
\label{ss:cogs_cbn}
             

The modification to obtain the concurrent version of OGS (\COGS) for call-by-name are
similar to those for call-by-value in Section~\ref{s:cogs}. 

Below are the rules for the call-by-name  OGS (\COGS).  
We use the same notations as the call-by-value version. 

 \setlength{\arraycolsep}{3pt}%
 \[
  \begin{array}{l|lll@{}l}
   (Int) & \conf{A \cdot \pmap{p}{M},\gamma,\phi} & \arrC{\qquad} & 
     \conf{A \cdot \pmap{p}{N},\gamma} & \text{ when } M \redn N\\
%
%
   (PA) & \conf{A \cdot \pmap{p}{V},\gamma,\phi} & \arrC{\ansP{p}{v}} & 
     \conf{A, \gamma \cdot \pmap{v}{V},\phi\uplus \{v\}} \\
%
   (PVQ) & \conf{A \cdot \pmap{p}{E[vM]},\gamma,\phi} & \arrC{\questPV{v}{y}{q}} & 
     \multicolumn{2}{l}{\conf{A, \gamma \cdot \pmap{y}{M}\cdot \pmap{q}{(E,p)},\phi\uplus \{y,q\}}} \\
%
   (PTQ) & \conf{A \cdot \pmap{p}{E[x]},\gamma,\phi} & \arrC{\questPT{x}{q}} & 
     \conf{A, \gamma \cdot \pmap{q}{(E,p)},\phi \uplus \{q\}} \\
%
   (OA) & \conf{A, \gamma,\phi} & \arrC{\ansO{p}{v}} & 
     \conf{A \cdot \pmap{q}{E[v]},\gamma,\phi \uplus \{v\}} & \text{ when } \gamma(p) = (E,q)  \\

%
   (OVQ) & \conf{A, \gamma \cdot [v \mapsto V],\phi} & \arrC{\questOV{v}{y}{p}} & 
     \conf{A \cdot \pmap{p}{V y},\gamma,\phi \uplus \{p,y\}} &\\
%
   (OTQ) & \conf{A,\gamma,\phi} & \arrC{\questOT{x}{p}} & 
     \conf{A \cdot \pmap{p}{M},\gamma,\phi \uplus \{p\}} & \text{ when } \gamma(x) = M\\
%
   (IOQ) &\initconf{\phi}{M}  & 
 \arrC{\questOinit{p}}
 &
     \conf{M,p,\emptymap,\phi \uplus \{p\}}
   % (IOQ) & \initconf{M}{\phi} & \arrC{\questOinit{p}} &
   %   \conf{M,p,\emptymap,\phi \cup \{p\}}
  \end{array}
 \]
\noindent
As usual, for encoding   \COGS,
it suffices  to consider
 running
terms and configurations,  as the encoding remains otherwise the same.
The encoding of the running term is: 
\[ 
\begin{array}{rcl}
\encoN{\pmap{p}{M}\cdot A  } & \defi & 
\encoNa  {M}{p} | \encoN{ A  } \\
\encoN{\emptymap} & \defi & \nil
\end{array} 
\]
The encoding of configurations is then:
\[
\encoN{\conf{ A , \gamma, \phi } }
\defi  \pairP{  \encoN A
 | \encoN \gamma }
\]



\noindent
We report the operational correspondence for strong transitions.

\begin{lemma}[from \COGS to \piI, strong transitions]
\label{l:CBNpogs_short}
$ $ \begin{enumerate}

\item  If $F \longrightarrowp F'$ then 
$ \encoSN F \longrightarrowPI \;  \contr \encoSN {F'}$; 

\item  If $F \arrC{\ell} F'$  then 
$ \encoSN F \arrPI{\ell}\;  \wbPI  \encoSN  {F'}$.
\end{enumerate}
\end{lemma}

 


\begin{lemma}[from  \piI to \COGS, strong transitions]
\label{l:CBNpogs_short_converse}
$ $ \begin{enumerate}

\item  If
$ \encoSN F \longrightarrowPI  P$, then there is $F'$ s.t.\  
 $F \longrightarrowp F'$ and 
$P 
\contr \encoSN {F'}$;

\item  If
$ \encoSN F \arrPI{\ell}  P$, then there is $F'$ s.t.\  
 $F \arrC{\ell} F'$ and 
$P 
\wbPI \encoSN {F'}$.
\end{enumerate}
\end{lemma} 



\begin{lemma} %[full abstraction for  traces and bisimilarity] %  ordinary LTS]
\label{l:ogsp_pi_traces_bisimulation_cbn}
For any $\FF,\FF'$ we have:
\begin{enumerate}
\item
 $\FF_1 \TEp \FF_2$ iff 
 $\encoN {\FF_1} \TEPI \encoN {\FF_2}$;  
\item
 $\FF_1 \wbC \FF_2$ iff 
 $\encoN {\FF_1} \wbPI  \encoN {\FF_2}$.  
\end{enumerate} 
\end{lemma} 


\begin{theorem}
\label{t:CBNfa_all}
For all $M,N$   with $\fv{M,N}\subseteq \phi$, we have: 
$\conf{ M,\phi} \TEp \conf{ N,\phi}$ iff 
$\conf{ M,\phi} \wbC \conf{ N,\phi}$ iff 
$\pairP{\encoSN M} \TEPI \pairP{\encoSN N}$
iff 
$\pairP{\encoSN M} \wb \pairP{\encoSN N}$. 
\end{theorem} 


\subsection{Up-to techniques for games, and relationship between Concurrent and Alternating OGS}
\label{ss:upto_games_cbn}

As for call-by-value, so for call-by-name the 
'up-to
  composition' technique for bisimulation 
(Definition~\ref{d:uptoComp}) 
can be imported from the $\pi$-calculus, and then used to establish that the equivalence
induced on $\lambda$-terms by the sequential and concurrent versions of OGS  coincide. 

We use  the technique to prove the call-by-name analogue of 
Lemma~\ref{l:OGS_COGS}. 



\begin{corollary}
\label{c:OGS_COGS_cbn}
For any   $\lambda$-terms $M,N$, the following statement are the same: 
$\conf{M} \TEA \conf {N}$;
$\conf{M} \wbA \conf {N}$;
$\conf{M} \TEp \conf {N}$;
$\conf{M} \wbC \conf {N}$;
$\encoN {{M} } \TEN\encoN {{N}}$; 
$\encoN {{M}} \wbOPI \encoN {{N}}$; 
$\encoN {{M} } \TEPI\encoN {{N}}$; 
$\encoN {{M}} \wbPI \encoN {{N}}$. 
\end{corollary} 
\noindent
The equivalence induced by the \piI (or  $\pi$-calculus) encoding of the call-by-name
$\lambda$-calculus is known to coincide with    that of the  
 L{\'e}vy-Longo Trees \cite{San95lazy}.  In the light of  
Corollary~\ref{c:OGS_COGS_cbn}, the result can thus be transported to OGS, both 
in its Alternating and its Concurrent variants, 
%\AOGS and \COGS, 
and both for traces and for bisimilarity.

Connections between game semantics and tree representation of $\lambda$-terms
have been previously explored, e.g.,  in~\cite{ong2004games}, 
where a game model of untyped call-by-name $\lambda$-calculus
is shown to correspond to Levy-Longo trees.

The connections that we have derived above, however, go beyond
L{\'e}vy-Longo Trees:
in particular,
the bisimilarities that have been established
between the observables (i.e., the dynamics) in the two models
set a tight relationship between them 
while allowing one to transfer results and techniques along the lines of what we have
shown for call-by-value.  
